% ---------------------------------------------------------------
% Currículo — Leandro Rocha de Brito | Desenvolvedor Mobile Flutter
% Compilar com: xelatex curriculo-leandro-rocha-flutter.tex
% ---------------------------------------------------------------
% Decisões ATS-safe (NÃO alterar sem revalidar a extração de texto):
%   - Coluna única, sem tabelas de layout, sem ícones
%   - Contato no corpo do documento, nunca em header/footer
%   - Ligatures=NoCommon evita que "fi"/"fl" quebrem a extração
%   - RawFeature={-calt;-case} desliga as alternativas contextuais da
%     Inter, que trocam "(", ")" e "-" por glifos sem mapeamento
%     ToUnicode — sem isso, parênteses e hífens somem na extração
%   - Hifenização desligada para não partir palavras-chave
% ---------------------------------------------------------------
% Decisões de legibilidade:
%   - Texto NÃO justificado (ragged right): justificar sem hifenização
%     produz espaçamento irregular e "rios" brancos entre palavras
%   - Corpo 9pt (classe extarticle), entrelinha 1.12, respiro entre blocos
%   - Uma única cor de destaque (petróleo) para criar hierarquia
% ---------------------------------------------------------------
\documentclass[9pt,a4paper]{extarticle}

\usepackage{fontspec}
\usepackage[margin=1.35cm,top=1.2cm,bottom=1.1cm]{geometry}
\usepackage{xcolor}
\usepackage[hidelinks]{hyperref}

% Metadados do PDF: nomeiam a aba do leitor e preenchem os campos que
% alguns ATS leem antes mesmo de extrair o texto. pdflang ajuda leitores
% de tela a escolher a pronúncia correta.
\hypersetup{
  pdftitle={Leandro Rocha de Brito — Desenvolvedor Mobile Flutter},
  pdfauthor={Leandro Rocha de Brito},
  pdfsubject={Currículo — Desenvolvedor Mobile Flutter},
  pdfkeywords={Flutter, Dart, Mobile, Clean Architecture, Riverpod, BLoC, ValueNotifier, ChangeNotifier, MVVM, Monorepo, Microapps, Testes, TDD, Firebase, Supabase},
  pdflang={pt-BR}
}
\usepackage{parskip}

% --- Paleta -----------------------------------------------------
% Todos os tons derivam de #120A8F (H=244°) variando apenas saturação
% e luminosidade. Entre parênteses, a luminância Rec.709 convertida em
% % de cinza — é o que uma impressora monocromática enxerga.
% Nenhum tom passa de 36%: acima disso o toner P&B apaga o traço.
\definecolor{ink}{HTML}{181820}      % texto principal        (10% cinza)
\definecolor{accent}{HTML}{120A8F}   % nome, seções, empresas ( 8% cinza)
\definecolor{rule}{HTML}{0000CD}     % filetes e marcadores   ( 6% cinza)
\definecolor{muted}{HTML}{000080}    % rótulos e metadados    ( 4% cinza)

% --- Fontes -----------------------------------------------------
\setmainfont{Inter}[
  Path           = /Users/leandrorocha/Library/Fonts/,
  UprightFont    = Inter-Regular.ttf,
  ItalicFont     = Inter-Italic.ttf,
  BoldFont       = Inter-SemiBold.ttf,
  BoldItalicFont = Inter-BoldItalic.ttf,
  Ligatures      = {NoCommon,NoContextual},
  RawFeature     = {-calt;-case},
]
\newfontfamily\interbold{Inter}[
  Path        = /Users/leandrorocha/Library/Fonts/,
  UprightFont = Inter-Bold.ttf,
  Ligatures   = {NoCommon,NoContextual},
  RawFeature  = {-calt;-case},
]
% Kerning desligado (-kern). Usar em trechos curtos cujo par de glifos
% recebe kern forte da Inter ("PJ", "7."): o xdvipdfmx grava esse ajuste
% como deslocamento no operador TJ e o extrator de texto do ATS o lê como
% espaço, produzindo "P J" e "17 .6K". Sem kern, a extração sai limpa.
\newfontfamily\nokern{Inter}[
  Path        = /Users/leandrorocha/Library/Fonts/,
  UprightFont = Inter-Regular.ttf,
  Ligatures   = {NoCommon,NoContextual},
  RawFeature  = {-calt;-case;-kern},
]
% Títulos de seção: bold com espacejamento entre letras.
\newfontfamily\secfont{Inter}[
  Path        = /Users/leandrorocha/Library/Fonts/,
  UprightFont = Inter-Bold.ttf,
  LetterSpace = 8.0,
  Ligatures   = {NoCommon,NoContextual},
  RawFeature  = {-calt;-case},
]

\color{ink}
\setlength{\parindent}{0pt}
\setlength{\parskip}{0.42em}
\linespread{1.12}
\pagestyle{empty}
\raggedright   % não justificar — ver nota de legibilidade acima

% Hifenização desligada: palavras partidas quebram o match de keyword.
\hyphenpenalty=10000
\exhyphenpenalty=10000

% --- Comandos ---------------------------------------------------
% Seção: caixa alta espaçada, cor de destaque, filete abaixo.
\newcommand{\sect}[1]{%
  \vspace{0.85em}%
  {\secfont\fontsize{9.5}{11}\selectfont\color{accent}\MakeUppercase{#1}}\par
  \vspace{0.30em}%
  {\color{rule}\hrule height 1.1pt}%
  \vspace{0.75em}%
}

% Experiência: empresa — cargo | vínculo | url ......... período
% A URL fica na MESMA linha do cabeçalho, aproveitando o espaço vazio antes
% da data: impresso o link clicável se perde, mas uma linha própria por
% experiência custaria uma página inteira. Corpo menor por hierarquia — é
% metadado de verificação, não conteúdo.
\newcommand{\job}[5]{%
  \vspace{0.75em}%
  {\interbold\fontsize{11}{13}\selectfont\color{accent}#1}%
  \hspace{0.35em}{\interbold\fontsize{10.5}{13}\selectfont — #2}%
  \hspace{0.35em}{\color{muted}\fontsize{9}{12}\selectfont | \mbox{#3}}%
  \hspace{0.45em}{\color{rule}\fontsize{7.8}{11}\selectfont #5}%
  \hfill{\color{muted}\fontsize{9}{12}\selectfont #4}\par
  \vspace{0.30em}%
}

% Rótulo de campo repetitivo (Descrição da Empresa, Projeto...) → discreto
\newcommand{\fld}[1]{{\color{muted}\textbf{#1}}}
% Rótulo dentro de um bullet (informativo) → cor do texto
\newcommand{\lbl}[1]{\textbf{#1}}

% Lista compacta com o ambiente `list` nativo (evita depender de enumitem).
\newenvironment{tight}{%
  \begin{list}{\textcolor{rule}{\rule[0.18em]{0.32em}{0.32em}}}{%
    \setlength{\leftmargin}{1.15em}%
    \setlength{\labelwidth}{0.60em}%
    \setlength{\labelsep}{0.55em}%
    \setlength{\topsep}{0.30em}%
    \setlength{\itemsep}{0.32em}%
    \setlength{\parsep}{0pt}%
    \setlength{\partopsep}{0pt}%
    \setlength{\rightmargin}{0pt}%
  }%
}{\end{list}}

\begin{document}

% =================== CABEÇALHO ===================
{\interbold\fontsize{21}{24}\selectfont\color{accent} LEANDRO ROCHA DE BRITO}\par
\vspace{0.20em}
{\fontsize{12}{15}\selectfont\color{ink} Desenvolvedor Mobile Flutter Sênior}\par
\vspace{0.60em}
{\fontsize{9.3}{13}\selectfont\color{muted}
\href{https://www.linkedin.com/in/leandrorochaadm}{https://linkedin.com/in/leandrorochaadm} \quad
\href{https://github.com/leandrorochaadm}{https://github.com/leandrorochaadm} \quad
\href{http://wa.me/5511997118447}{(11) 99711-8447}\par
\vspace{0.15em}
\href{mailto:leandrorochaadm@gmail.com}{leandrorochaadm@gmail.com} \quad
100\% remoto\par
}
\vspace{0.55em}
{\color{rule}\hrule height 1.1pt}

% =================== APRESENTAÇÃO ===================
\sect{Apresentação}

Desenvolvedor mobile com \textbf{mais de 4 anos dedicados a Flutter,} construindo aplicativos mobile para e-commerce, saúde e fintech, incluindo projetos para a \textbf{Claro}. Trabalho com \textbf{Clean Architecture} e \textbf{testes unitários, de widget e de integração}, com apps publicados na \textbf{Google Play} e na \textbf{App Store}. Atuo em \textbf{decisões arquiteturais} e na \textbf{definição dos padrões de qualidade e testes} adotados pelo time.

% =================== HABILIDADES ===================
\sect{Habilidades}

\begin{tight}
  \item \lbl{Flutter e Dart:} \textbf{BLoC} | \textbf{Riverpod} | \textbf{ValueNotifier} | \textbf{ChangeNotifier} | Provider | Get It | Dio | DartZ | Hive | Flutter Secure Storage | fl\_chart | Google Maps | Flutter Web (PWA) | Design System próprio | Freezed | GoRouter | build\_runner
  \item \lbl{Plataformas e recursos nativos:} \textbf{Android} | \textbf{iOS} | \textbf{Biometria facial (FaceTec, ML Kit)} | \textbf{integração de SDKs e plugins nativos} (Kotlin/Swift)
  \item \lbl{Arquitetura:} \textbf{Clean Architecture} | MVVM | MVC | \textbf{tratamento de erro tipado (Either/Result + Failure)} | \textbf{Monorepo (Dart Workspace)} | \textbf{modularização em packages} | microapps | MultiRepo | modelagem multi-tenant
  \item \lbl{IA no desenvolvimento:} uso de \textbf{IA como apoio ao desenvolvimento} com \textbf{regras de projeto versionadas} para garantir aderência ao padrão do time, e \textbf{validação técnica em code review}
  \item \lbl{Qualidade e observabilidade:} \textbf{Testes unitários e de widget (Mocktail)} | \textbf{TDD} | \textbf{Datadog} | \textbf{Firebase Crashlytics} | Microsoft Clarity | \textbf{CI/CD (GitHub Actions, Codemagic)} | \textbf{Git Flow} e \textbf{trunk-based} | code review
  \item \lbl{Firebase:} Analytics | Firestore | Storage | Remote Config | Cloud Messaging (FCM)
  \item \lbl{Backend e dados:} API REST | \textbf{Supabase / PostgreSQL} (Row Level Security, migrations versionadas) | SQLite | MySQL | Cloudflare Workers
  \item \lbl{Distribuição e localização:} publicação e manutenção de apps na \textbf{Google Play} e \textbf{App Store} | push notifications (FCM e OneSignal) | Internacionalização (i18n) e localização (l10n)
  \item \lbl{Outras:} \textbf{Git} | GitHub | GitLab | \textbf{Scrum} | Kanban | Jira | Figma | Linux | macOS | VueJS | Astro | Golang
\end{tight}

% =================== EXPERIÊNCIAS ===================
\sect{Experiências}

\job{\href{https://clyvo.global/}{Clyvo}}{Desenvolvedor Flutter}{{\nokern PJ}}{Outubro/2025 até Julho/2026}{\href{https://clyvo.global/}{https://clyvo.global}}


\fld{Descrição da Empresa:} Healthtech brasileira de telemedicina, que conecta médicos e pacientes em consultas online.

\fld{Projeto:} Aplicativo de telemedicina onde o paciente busca o médico por especialidade, agenda a consulta e é atendido por videochamada, recebendo a receita assinada digitalmente ao fim do atendimento. Conta ainda com sala de espera virtual, chat entre médico e paciente, prontuário, prescrição, assistente de IA por voz, cadastro com biometria facial e notificações push.

\fld{Responsabilidades Técnicas:}
\begin{tight}
  \item \lbl{Qualidade e testes:} estruturei \textbf{do zero} a estratégia e os padrões de teste adotados pelo time, levando a base a mais de \textbf{6.800 testes} unitários e de widget em Mocktail e \textbf{85\% de cobertura} do código; \textbf{introduzi a prática de code review} no time mobile, atuando como revisor. \textbf{POC de testes E2E com Patrol} e comparativo com Appium para o QA.
  \item \lbl{Arquitetura:} \textbf{conduzi a migração de multirepo para monorepo} em Dart Workspace, estruturando os \textbf{31 módulos de negócio} com Clean Architecture. Gerência de estado com \textbf{Riverpod} no app e \textbf{BLoC} no SDK de chat com IA. \textbf{Design system e SDK de chat isolados em packages próprios,} compartilhados entre os apps do monorepo.
  \item \lbl{Teleconsulta em tempo real:} manutenção e evolução do módulo de videochamada em LiveKit/WebRTC, com tela de chamada nativa (CallKit no iOS e IncomingCall no Android).
  \item \lbl{Onboarding e documentos:} \textbf{integração da biometria facial} (FaceTec e ML Kit) no fluxo de cadastro, com tratamento de falha e cancelamento de captura; tratamento de erro e reautenticação no fluxo de \textbf{assinatura digital}.
  \item \lbl{Integrações e infraestrutura mobile:} consumo de APIs REST com Dio, persistência com Hive e Secure Storage, observabilidade com \textbf{Datadog}.
\end{tight}

\job{\href{https://www.monetizze.com.br/}{Monetizze}}{Desenvolvedor Flutter}{{\nokern PJ}}{Janeiro/2025 até Setembro/2025}{\href{https://www.monetizze.com.br/}{https://monetizze.com.br}}

\fld{Descrição da Empresa:} Plataforma de vendas online para produtores e afiliados de produtos digitais e físicos.

\fld{Projeto:} Aplicativo mobile para afiliados e produtores digitais, com gestão de vendas, comissões e saques, relatórios, métricas de desempenho e notificações em tempo real. \textbf{1M+ downloads} · nota 4,8/5 na Google Play.

\fld{Responsabilidades Técnicas:} Desenvolvimento em Flutter com \textbf{Clean Architecture e BLoC}; APIs REST; Firebase (Crashlytics, Analytics, Remote Config); SQLite e Secure Storage; autenticação biométrica; push via OneSignal e FCM; i18n; gráficos interativos; testes unitários e code review. Build e publicação em \textbf{Codemagic} com GitLab.

\job{\href{https://sof.to/}{SOFTO}}{Desenvolvedor Flutter}{{\nokern PJ}}{Outubro/2023 até Novembro/2024}{\href{https://sof.to/}{https://sof.to} · \href{https://www.thefaapp.org}{https://thefaapp.org}}

\fld{Descrição da Empresa:} Empresa de tecnologia com sedes no Brasil e nos EUA, atendendo clientes como OMS, FGV e Stone.

\fld{Projeto \href{https://www.thefaapp.org}{FA app}:} Aplicativo global que conecta médicos, cientistas e famílias na luta contra a Ataxia de Friedreich, doença neurodegenerativa rara. \textbf{10K+ downloads} · nota 4,6/5.

\fld{Responsabilidades Técnicas:} Refatoração visual e evolução do app 100\% Flutter, com \textbf{MVVM e gerência de estado em ChangeNotifier}, melhorando usabilidade e acessibilidade. Testes unitários e de widget, Firebase Analytics e Firestore, internacionalização (i18n), leitura de artigos e notícias e compartilhamento via chat interno e redes sociais.

\job{\href{https://amaris.com/}{Amaris Consulting}}{Desenvolvedor Flutter}{{\nokern PJ}}{Abril/2022 até Agosto/2023}{\href{https://amaris.com/}{https://amaris.com} · \href{https://claropay.com.br/}{https://claropay.com.br}}


\fld{Descrição da Empresa:} Atuei em um projeto chave para a \href{https://claropay.com.br/}{Claro Pay}, fintech do grupo Claro focada em serviços financeiros digitais.

\fld{Projeto seguros e assistências técnicas:} Responsável pelo desenvolvimento end-to-end do módulo de seguros e assistências técnicas, com funcionalidades de contratação, comparação de produtos, gestão de contratos ativos (cancelamento e alteração) e suporte ao cliente. \textbf{1M+ downloads} na Google Play.

\fld{Responsabilidades Técnicas:} Apliquei Clean Architecture e boas práticas de código limpo, desenvolvendo o módulo dentro da \textbf{arquitetura de microapps} adotada no projeto, com \textbf{monorepo e módulos isolados em packages} e gerência de estado em \textbf{ValueNotifier.} Utilizei TDD com testes unitários e de widget para garantir robustez, integração com Firebase e APIs REST.

% =================== PROJETOS PESSOAIS ===================
\sect{Projetos Pessoais e Voluntariado}

\job{\href{http://lions-pontos.tektonsoftwares.workers.dev}{Lions Pontos}}{Lions Club de Ji-Paraná Centro}{Voluntário}{Julho/2026 até o momento}{}

\fld{Projeto:} Plataforma multi-clube que apura em tempo real o prêmio de associado mais atuante de clubes de serviço, substituindo a apuração manual em atas de papel por um ranking público e auditável. Lançamento de presença pelo celular, extrato individual por associado, tabela de pontos configurável por clube e exportação do ano leonístico. Em piloto no Lions Clube de Ji-Paraná/RO.

\fld{Demo:} \href{http://lions-pontos.tektonsoftwares.workers.dev/demo}{https://lions-pontos.tektonsoftwares.workers.dev/demo}\\
\fld{Código:} \href{http://github.com/leandrorochaadm/lions-club-points}{https://github.com/leandrorochaadm/lions-club-points}

\fld{Responsabilidades Técnicas:} \textbf{Único responsável por todas as decisões técnicas do produto — do modelo de dados ao deploy.}
\begin{tight}
  \item \lbl{Frontend:} concepção e desenvolvimento end-to-end em \textbf{Flutter Web} (PWA instalável) com Clean Architecture e Riverpod.
  \item \lbl{Backend e dados:} \textbf{Supabase/Postgres} com modelagem \textbf{multi-tenant} e isolamento entre clubes garantido por \textbf{Row Level Security no banco} (não na aplicação); \textbf{19 migrations} versionadas.
  \item \lbl{Qualidade:} \textbf{968 testes} automatizados com mais de \textbf{90\% de cobertura}, com \textbf{CI no GitHub Actions barrando merge} abaixo de 80\%.
  \item \lbl{Infraestrutura:} deploy em \textbf{Cloudflare Workers}.
\end{tight}

% =================== FORMAÇÃO ===================
\sect{Formação Acadêmica}

{\interbold\fontsize{10.5}{13}\selectfont\color{accent} MBA em Inteligência Artificial e Big Data}\hspace{0.35em}{\interbold\fontsize{10.5}{13}\selectfont — \href{https://mba.iabigdata.icmc.usp.br/}{ICMC/USP}}\par
{\color{muted}\fontsize{8.6}{11}\selectfont Julho/2025 até Outubro/2026, em andamento \quad \href{https://mba.iabigdata.icmc.usp.br/}{https://mba.iabigdata.icmc.usp.br}}\par
\vspace{0.20em}
Programa de 400h em machine learning, NLP e engenharia de dados em cloud, com projeto aplicado.

\vspace{0.60em}
{\interbold\fontsize{10.5}{13}\selectfont\color{accent} Análise e Desenvolvimento de Sistemas}\hspace{0.35em}{\interbold\fontsize{10.5}{13}\selectfont — IFRO} {\color{muted}\fontsize{9}{12}\selectfont (Instituto Federal de Rondônia)}\par
{\color{muted}\fontsize{8.6}{11}\selectfont Fevereiro/2020 até Novembro/2022}\par

\sect{Idiomas}

\textbf{Português} — nativo \quad · \quad \textbf{Inglês técnico} — leitura fluente de documentação, código e especificações

\end{document}
